% GENERATED FILE. Edit canonical lessons, then run npm run generate:books.
% canonical-source-hash: fnv1a64:ae38a0b48a39e7ee
% canonical-lessons: BN-C408-basa, BN-C408-bharate, BN-C408-ogrim, BN-C408-jamanot, BN-C408-bosbas-kora

\chapter{Rented Rooms}
\label{ch:bn-rented-rooms}
\hlchaptermodality{408}

\begin{chapteropening}
\textbf{By the end of this chapter:} \emph{I can rent a flat, pay an advance and a deposit, and say where I live.}

Complete the last lesson of chapter 408: I can rent a flat, pay an advance and a deposit, and say where I live.
\end{chapteropening}

\section[\texorpdfstring{\bn{বাসা} (bāsā)}{বাসা (bāsā)}]{\bn{বাসা} (bāsā) — a rented home; lodgings; a nest}

\label{lesson:BN-C408-basa}



\begin{quote}
Before the new one: say the Bengali for a tip (for service), then the Bengali for a laundry.
\end{quote}

\subsection*{You'll want to know: \bn{বাসা}}
\textbf{\bn{বাসা}} — \emph{bāsā} — \textquotedblleft{}a rented home; lodgings; a nest\textquotedblright{}.

\begin{quote}
\bn{শহরে} \bn{মীরা} \bn{একটা} \bn{ছোট} \bn{বাসা} \bn{ভাড়া} \bn{নিল।}
\end{quote}

\begin{quote}
\emph{shôhôre mīrā ekṭā chhoṭo bāsā bhāṛā nilô.}
\end{quote}

\begin{quote}
In the city Mira rented a small flat.
\end{quote}

Also a bird's nest: \bn{পাখির} \bn{বাসা}. \bn{বাড়ি} is the family home; \bn{বাসা} is where you live now.

\begin{cousinweb}
\textbf{\bn{বাসা}} (\emph{bāsā}) is inherited: CDIAL lists Bengali \emph{bāsā}, \textquotedblleft{}a lodging-house, a nest, a lair\textquotedblright{}, under Sanskrit \emph{vāsá}, \textquotedblleft{}a dwelling\textquotedblright{} (CDIAL 11591).
\end{cousinweb}

\begin{grammarlens}[title={What you've built}]
One more word for getting around.
\end{grammarlens}

\subsection*{Your turn}
\begin{itemize}
\raggedright
  \item \emph{Say it:} \emph{bāsā}
  \item \emph{Say it:} \emph{bāsā}, once more
  \item \emph{Say it:} \emph{shôhôre mīrā ekṭā chhoṭo bāsā bhāṛā nilô.} — In the city Mira rented a small flat.
  \item \emph{Recall:} say the Bengali for a tip (for service), then the Bengali for a laundry, then say \emph{bāsā} again
\end{itemize}

\begin{tcolorbox}[breakable,colback=teal!4,colframe=teal!35!black,title={Before you move on}]
What does \bn{বাসা} mean? (\textquotedblleft{}A rented home; lodgings; a nest\textquotedblright{}.) Say it once more.
\end{tcolorbox}
\section[\texorpdfstring{\bn{ভাড়াটে} (bhāṛāṭe)}{ভাড়াটে (bhāṛāṭe)}]{\bn{ভাড়াটে} (bhāṛāṭe) — a tenant}

\label{lesson:BN-C408-bharate}



\begin{quote}
Before the new one: say the Bengali for a laundry, then the Bengali for a rented home.
\end{quote}

\subsection*{You'll want to know: \bn{ভাড়াটে}}
\textbf{\bn{ভাড়াটে}} — \emph{bhāṛāṭe} — \textquotedblleft{}a tenant\textquotedblright{}.

\begin{quote}
\bn{আমাদের} \bn{বাড়িতে} \bn{নতুন} \bn{ভাড়াটে} \bn{এল।}
\end{quote}

\begin{quote}
\emph{āmāder bāṛite nôtun bhāṛāṭe elô.}
\end{quote}

\begin{quote}
A new tenant moved into our house.
\end{quote}

\begin{cousinweb}
\textbf{\bn{ভাড়াটে}} (\emph{bhāṛāṭe}) is built on \textbf{\bn{ভাড়া}}, \textquotedblleft{}rent\textquotedblright{}, which CDIAL lists as Bengali \emph{bhāṛā}, \textquotedblleft{}hire, rent\textquotedblright{}, under the reconstructed \emph{bhārta} (CDIAL 9468), with the describing ending -\bn{টে}: of rent, hired; as a noun it settled on the one who pays the rent.
\end{cousinweb}

\begin{grammarlens}[title={What you've built}]
One more word for getting around.
\end{grammarlens}

\subsection*{Your turn}
\begin{itemize}
\raggedright
  \item \emph{Say it:} \emph{bhāṛāṭe}
  \item \emph{Say it:} \emph{bhāṛāṭe}, once more
  \item \emph{Say it:} \emph{āmāder bāṛite nôtun bhāṛāṭe elô.} — A new tenant moved into our house.
  \item \emph{Recall:} say the Bengali for a laundry, then the Bengali for a rented home, then say \emph{bhāṛāṭe} again
\end{itemize}

\begin{tcolorbox}[breakable,colback=teal!4,colframe=teal!35!black,title={Before you move on}]
What does \bn{ভাড়াটে} mean? (\textquotedblleft{}A tenant\textquotedblright{}.) Say it once more.
\end{tcolorbox}
\section[\texorpdfstring{\bn{অগ্রিম} (ôgrim)}{অগ্রিম (ôgrim)}]{\bn{অগ্রিম} (ôgrim) — an advance (payment); in advance}

\label{lesson:BN-C408-ogrim}



\begin{quote}
Before the new one: say the Bengali for a rented home, then the Bengali for a tenant.
\end{quote}

\subsection*{You'll want to know: \bn{অগ্রিম}}
\textbf{\bn{অগ্রিম}} — \emph{ôgrim} — \textquotedblleft{}an advance (payment); in advance\textquotedblright{}.

\begin{quote}
\bn{বাসার} \bn{জন্য} \bn{দুই} \bn{মাসের} \bn{ভাড়া} \bn{অগ্রিম} \bn{দিতে} \bn{হবে।}
\end{quote}

\begin{quote}
\emph{bāsār jônyô dui māser bhāṛā ôgrim dite hôbe.}
\end{quote}

\begin{quote}
For the flat you have to pay two months' rent in advance.
\end{quote}

\begin{cousinweb}
\textbf{\bn{অগ্রিম}} (\emph{ôgrim}) is Sanskrit \emph{agrima}, \textquotedblleft{}being in front, preceding\textquotedblright{} (Monier-Williams), from \emph{agra}, the front: money paid before.
\end{cousinweb}

\begin{grammarlens}[title={What you've built}]
One more word for getting around.
\end{grammarlens}

\subsection*{Your turn}
\begin{itemize}
\raggedright
  \item \emph{Say it:} \emph{ôgrim}
  \item \emph{Say it:} \emph{ôgrim}, once more
  \item \emph{Say it:} \emph{bāsār jônyô dui māser bhāṛā ôgrim dite hôbe.} — For the flat you have to pay two months' rent in advance.
  \item \emph{Recall:} say the Bengali for a rented home, then the Bengali for a tenant, then say \emph{ôgrim} again
\end{itemize}

\begin{tcolorbox}[breakable,colback=teal!4,colframe=teal!35!black,title={Before you move on}]
What does \bn{অগ্রিম} mean? (\textquotedblleft{}An advance (payment); in advance\textquotedblright{}.) Say it once more.
\end{tcolorbox}
\section[\texorpdfstring{\bn{জামানত} (jāmānôt)}{জামানত (jāmānôt)}]{\bn{জামানত} (jāmānôt) — a security deposit}

\label{lesson:BN-C408-jamanot}



\begin{quote}
Before the new one: say the Bengali for a tenant, then the Bengali for an advance (payment).
\end{quote}

\subsection*{You'll want to know: \bn{জামানত}}
\textbf{\bn{জামানত}} — \emph{jāmānôt} — \textquotedblleft{}a security deposit\textquotedblright{}.

\begin{quote}
\bn{ঘর} \bn{ছাড়ার} \bn{সময়} \bn{মালিক} \bn{জামানতের} \bn{টাকা} \bn{ফিরিয়ে} \bn{দিলেন।}
\end{quote}

\begin{quote}
\emph{ghôr chhāṛār sômôy mālik jāmānôter ṭākā phiriye dilen.}
\end{quote}

\begin{quote}
When we left the room, the owner gave back the deposit.
\end{quote}

\begin{cousinweb}
\textbf{\bn{জামানত}} (\emph{jāmānôt}) is usually traced to Persian \emph{zamānat}, \textquotedblleft{}surety, a guarantee\textquotedblright{}, from Arabic \emph{ḍamāna}.
\end{cousinweb}

\begin{grammarlens}[title={What you've built}]
One more word for getting around.
\end{grammarlens}

\subsection*{Your turn}
\begin{itemize}
\raggedright
  \item \emph{Say it:} \emph{jāmānôt}
  \item \emph{Say it:} \emph{jāmānôt}, once more
  \item \emph{Say it:} \emph{ghôr chhāṛār sômôy mālik jāmānôter ṭākā phiriye dilen.} — When we left the room, the owner gave back the deposit.
  \item \emph{Recall:} say the Bengali for a tenant, then the Bengali for an advance (payment), then say \emph{jāmānôt} again
\end{itemize}

\begin{tcolorbox}[breakable,colback=teal!4,colframe=teal!35!black,title={Before you move on}]
What does \bn{জামানত} mean? (\textquotedblleft{}A security deposit\textquotedblright{}.) Say it once more.
\end{tcolorbox}
\section[\texorpdfstring{\bn{বসবাস} \bn{করা} (bôsbās kôrā)}{বসবাস করা (bôsbās kôrā)}]{\bn{বসবাস} \bn{করা} (bôsbās kôrā) — to live (in a place); to reside}

\label{lesson:BN-C408-bosbas-kora}



\begin{quote}
Before the new one: say the Bengali for an advance (payment), then the Bengali for a security deposit.
\end{quote}

\subsection*{You'll want to know: \bn{বসবাস} \bn{করা}}
\textbf{\bn{বসবাস} \bn{করা}} — \emph{bôsbās kôrā} — \textquotedblleft{}to live (in a place); to reside\textquotedblright{}.

\begin{quote}
\bn{এই} \bn{পাড়ায়} \bn{অনেক} \bn{পরিবার} \bn{বসবাস} \bn{করে।}
\end{quote}

\begin{quote}
\emph{ei pāṛāy ônek pôribār bôsbās kôre.}
\end{quote}

\begin{quote}
Many families live in this neighbourhood.
\end{quote}

Past (he/she …): \bn{তিনি} \bn{গ্রামে} \bn{বসবাস} \bn{করতেন}, \emph{tini grāme bôsbās kôrten}.

A careful word for settled living; \bn{থাকা} is the everyday one.

\begin{cousinweb}
\textbf{\bn{বসবাস} \bn{করা}} (\emph{bôsbās kôrā}) joins \bn{বস}-, the stem of \bn{বসা}, to sit, and \bn{বাস}, \textquotedblleft{}a dwelling\textquotedblright{}, Sanskrit \emph{vāsa} (Monier-Williams), to \bn{করা}, to do: to sit down and dwell.
\end{cousinweb}

\begin{grammarlens}[title={What you've built}]
One more word for getting around.
\end{grammarlens}

\subsection*{Your turn}
\begin{itemize}
\raggedright
  \item \emph{Say it:} \emph{bôsbās kôrā}
  \item \emph{Say it:} \emph{bôsbās kôrā}, once more
  \item \emph{Say it:} \emph{ei pāṛāy ônek pôribār bôsbās kôre.} — Many families live in this neighbourhood.
  \item \emph{Recall:} say the Bengali for an advance (payment), then the Bengali for a security deposit, then say \emph{bôsbās kôrā} again
\end{itemize}

\begin{tcolorbox}[breakable,colback=teal!4,colframe=teal!35!black,title={Before you move on}]
What does \bn{বসবাস} \bn{করা} mean? (\textquotedblleft{}To live (in a place); to reside\textquotedblright{}.) Say it once more.
\end{tcolorbox}
